% !TEX root = ../main.tex
% ============================================================
%  Fiches PSE1 — Bilan complémentaire · v0.2
%  Fichier : inc/commands.tex
%  Rôle    : Commandes maison utilisées par les sections.
%            Chaque commande possède deux variantes :
%              - suffixe A : version A3 (police standard, vspace généreux)
%              - suffixe B : version A4 (compactage, scriptsize, vspace serré)
%            Le dispatch en bas du fichier choisit la bonne variante via
%            \PaperFormat ; \ifA{...}{...} permet aussi des conditionnels
%            inline dans les fichiers de section.
% ============================================================

% ---- Hauteur fixe d'un champ \acronymfield (2 lignes pointillées) ----
\newlength{\fieldheight}

% ---- Police "intermédiaire" (XABCD, surveillance, gestes, glasgow) ----
\newcommand{\smallscript}{\fontsize{8.5}{10}\selectfont}  % defaut A3

% ---- Réglages spécifiques par format ----
\ifx\PaperFormat\PaperFormatA
  % A3
  \setlength{\fieldheight}{9.5mm}
\else
  % A4 : compactage agressif
  \setlength{\fieldheight}{6mm}
  \renewcommand{\smallscript}{\fontsize{6}{7}\selectfont}
\fi

% ============================================================
%  CHAMPS D'ACRONYME
%    #1 = couleur de section, #2 = lettre majuscule, #3 = filigrane
%    après la lettre, #4 = texte de référence (gris) ou options
% ============================================================

% ---- Champ 1 ligne : Lettre + filigrane + dotfill + ref grise à droite ----
\newcommand{\acronymfieldsingleA}[4]{%
  \par\noindent
  \textbf{\textcolor{#1}{\large #2}}%
  \rlap{\textcolor{#1!65}{\normalsize #3}}%
  {\color{pse-ligne}\dotfill}%
  {\textcolor{pse-gris!40}{\normalsize #4}}%
  \par\vspace{0.5mm}%
}
\newcommand{\acronymfieldsingleB}[4]{%
  \par\noindent
  \textbf{\textcolor{#1}{\large #2}}%
  \rlap{\textcolor{#1!65}{\small #3}}%
  {\color{pse-ligne}\dotfill}%
  {\textcolor{pse-gris!40}{\small #4}}%
  \par\vspace{0.2mm}%
}

% ---- Champ 2 lignes : 1 ligne titre + 1 ligne pointillée pour rédaction ----
\newcommand{\acronymfieldA}[4]{%
  \par\noindent
  \begin{minipage}[c][\fieldheight][t]{\linewidth}%
    \textbf{\textcolor{#1}{\large #2}}%
    \rlap{\textcolor{#1!65}{\normalsize #3}}%
    {\color{pse-ligne}\dotfill}\par\nobreak
    \noindent
    {\color{pse-ligne}\dotfill}%
    {\textcolor{pse-gris!40}{\normalsize #4}}%
  \end{minipage}%
  \par\vspace{0.5mm}%
}
\newcommand{\acronymfieldB}[4]{%
  \par\noindent
  \begin{minipage}[c][\fieldheight][t]{\linewidth}%
    \textbf{\textcolor{#1}{\large #2}}%
    \rlap{\textcolor{#1!65}{\small #3}}%
    {\color{pse-ligne}\dotfill}\par\nobreak
    \noindent
    {\color{pse-ligne}\dotfill}%
    {\textcolor{pse-gris!40}{\small #4}}%
  \end{minipage}%
  \par\vspace{0.2mm}%
}

% ---- Champ 1 ligne avec options à cocher (au lieu de la ref grise) ----
\newcommand{\acronymfieldcheckA}[4]{%
  \par\noindent
  \textbf{\textcolor{#1}{\large #2}}%
  \rlap{\textcolor{#1!65}{\normalsize #3}}%
  {\color{pse-ligne}\dotfill}\;#4%
  \par\vspace{0.5mm}%
}
\newcommand{\acronymfieldcheckB}[4]{%
  \par\noindent
  \textbf{\textcolor{#1}{\large #2}}%
  \rlap{\textcolor{#1!65}{\small #3}}%
  {\color{pse-ligne}\dotfill}\;#4%
  \par\vspace{0.2mm}%
}

% ---- Option à cocher (utilisée dans #4 de \acronymfieldcheck) ----
\newcommand{\optA}[1]{$\square$\,{\normalsize #1}\enspace}
\newcommand{\optB}[1]{$\square$\,{\small #1}\enspace}

% ---- Ligne libre : label en gras + pointillés (ex. Circonstances) ----
\newcommand{\writelineA}[1]{%
  \noindent\textbf{\normalsize #1}\enspace{\color{pse-ligne}\dotfill}\par\vspace{2.5pt}%
}
\newcommand{\writelineB}[1]{%
  \noindent\textbf{\small #1}\enspace{\color{pse-ligne}\dotfill}\par\vspace{0.5pt}%
}

% ---- Sous-titre de bloc (étiquette colorée pleine, ex. PQRST, FAST...) ----
\newcommand{\blocktitleA}[2]{%
  \vspace{1.5mm}%
  \noindent{\colorbox{#1}{\color{white}\bfseries\scriptsize\,#2\,}}%
  \vspace{1.5mm}\par
}
\newcommand{\blocktitleB}[2]{%
  \vspace{0.4mm}%
  \noindent{\colorbox{#1}{\color{white}\bfseries\scriptsize\,#2\,}}%
  \vspace{0.4mm}\par
}

% ---- Rond numéroté pour l'échelle de sévérité PQRST (0 à 10) ----
\newcommand{\sevA}[1]{%
  \tikz[baseline=-0.5ex]\node[draw=pse-jaune!50, circle, minimum size=5mm,
    inner sep=0pt, font=\smallscript]{#1};\!%
}
\newcommand{\sevB}[1]{%
  \tikz[baseline=-0.5ex]\node[draw=pse-jaune!50, circle, minimum size=3.5mm,
    inner sep=0pt, font=\smallscript]{#1};\!%
}

% ============================================================
%  DISPATCH : associe les noms publics (\acronymfield, \opt, ...) à la
%  variante A ou B selon le format de papier sélectionné.
% ============================================================
\ifx\PaperFormat\PaperFormatA
  \let\acronymfieldsingle\acronymfieldsingleA
  \let\acronymfield\acronymfieldA
  \let\acronymfieldcheck\acronymfieldcheckA
  \let\opt\optA
  \let\writeline\writelineA
  \let\blocktitle\blocktitleA
  \let\sev\sevA
\else
  \let\acronymfieldsingle\acronymfieldsingleB
  \let\acronymfield\acronymfieldB
  \let\acronymfieldcheck\acronymfieldcheckB
  \let\opt\optB
  \let\writeline\writelineB
  \let\blocktitle\blocktitleB
  \let\sev\sevB
\fi

% ---- Conditionnel inline pour les sections : \ifA{...A3...}{...A4...} ----
\newcommand{\ifA}[2]{%
  \ifx\PaperFormat\PaperFormatA #1\else #2\fi
}
